% ===== 底本 第 1 章《引言》；片段 tr01b；源块 B024--B037 =====

% [源 B024]
\begin{lemma*}
$\Hom^{\ast}(\mathcal{O}_p,\mathcal{O}_p)\cong\Lambda^{\ast}\langle x_1,\dotsb,x_n\rangle$ 是分次外代数，其 $n=\dim_{\mathbb{C}}\check{X}$ 个生成元的次数为 $1$。
\end{lemma*}

% [Tutor Note: B024 的 $n=\dim_{\mathbb{C}}\check{X}$ 依据：光滑闭点 $p$ 处的点对象 $\mathcal{O}_p$ 的余维数等于 $\dim\check{X}$；局部地 $\check{X}=\mathbb{C}^n$、$p=0$，正则参数 $x_1,\dotsb,x_n$ 共 $n$ 个，故 $\Lambda^{\ast}\langle x_1,\dotsb,x_n\rangle$ 恰有 $n$ 个次数 $1$ 的生成元。取 $n=1$ 可核验：$\operatorname{Ext}^{\ast}_{\mathbb{C}[x]}(\mathbb{C},\mathbb{C})\cong\Lambda^{\ast}\langle x\rangle$。]

% [源 B025]
\begin{proof}
这是一个局部问题，故可归约到 $\check{X}=\mathbb{C}^n$、$p=0$ 的情形，此时任务是用 $n$ 元多项式代数给出 $\mathbb{C}$ 的一个消解。这只需取\emph{Koszul 复形}
\begin{multline*}
  0\to \mathbb{C}[x_1,\dotsb,x_n]\otimes\Lambda^n\langle e_1,\dotsb,e_n\rangle
  \xrightarrow{d_n}\dotsb\xrightarrow{d_2}
  \mathbb{C}[x_1,\dotsb,x_n]\otimes\Lambda^1\langle e_1,\dotsb,e_n\rangle\\
  \xrightarrow{d_1}\mathbb{C}[x_1,\dotsb,x_n]\to 0,
\end{multline*}
其中微分由 $d_k(f\otimes e_{i_1}\wedge\dotsb\wedge e_{i_k})=\sum_{j=1}^k(-1)^{j-1}x_{i_j}f\otimes e_{i_1}\wedge\dotsb\wedge\widehat{e_{i_j}}\wedge\dotsb\wedge e_{i_k}$ 定义。
\end{proof}

% [Tutor Note: B025 的 Koszul 复形取 $n=1$ 退化核验：$0\to\mathbb{C}[x]\otimes\Lambda^1\langle e\rangle\xrightarrow{d_1}\mathbb{C}[x]\to0$ 即 $0\to\mathbb{C}[x]\xrightarrow{\cdot x}\mathbb{C}[x]\to0$，仅在 $H^0$ 处得到 $\mathbb{C}\cong\mathcal{O}_0$；微分中的符号 $(-1)^{j-1}$ 是保证 $d^2=0$ 的标准约定，$\widehat{e_{i_j}}$ 表示删去该因子。]

% [源 B026]
特别地，对象 $L\in\mathcal{Fuk}(X,\omega)$ 只有当 $L$ 的自同态代数等价于外代数 $\Lambda\langle x_1,\dotsb,x_n\rangle$ 时，才表现得像点对象。回忆对 $(X,\omega)$ 的两个拉格朗日子流形 $L_1,L_2$，$\mathcal{Fuk}(X,\omega)$ 中 $L_1$ 与 $L_2$ 之间的态射 $\Hom_{\mathcal{Fuk}(X,\omega)}(L_1,L_2)$ 由\emph{拉格朗日 Floer 复形} $CF^{\ast}(L_1,L_2)$ 给出：当 $L_1$ 与 $L_2$ 横截相交时，该复形是分次的，且由 $L_1$ 与 $L_2$ 的交点生成（一般情形下取相应的导出版本，参见 [FOOO09]）。当 $L_1$ 与 $L_2$ 是良好的闭拉格朗日子流形、$(X,\omega)$ 是良好的辛流形时，著名的 Floer 定理成立：

% [Tutor Note: B026 两侧的对照是本段关键：Floer 侧有 $HF^{\ast}(L,L;\mathbb{R})\cong\Lambda_{\mathbb{R}}\langle x_1,\dotsb,x_n\rangle$（见 B031），代数侧有 $\Hom^{\ast}(\mathcal{O}_p,\mathcal{O}_p)\cong\Lambda^{\ast}\langle x_1,\dotsb,x_n\rangle$（B024），都是“$n$ 个次数 $1$ 生成元的外代数”；易错点是把 $\Lambda^{\ast}$ 误读成对称代数 $\operatorname{Sym}^{\ast}$，而 B024 明确说它是分次外代数。]

% [源 B027]
\begin{theorem*}[{Floer'89 [Flo89]、Oh'96 [Oh96] 等}]
假设下列条件之一成立：
\begin{enumerate}[label=\roman*.,leftmargin=1.9em,itemsep=1pt,topsep=2pt]
% [源 B028]
  \item $X$ 是辛非球面的，且 $L\subseteq X$ 是闭的自旋拉格朗日子流形，满足 $\pi_2(X,L)=0$；
% [源 B029]
  \item $X$ 是单调辛流形，且 $L\subseteq X$ 是闭的单调自旋拉格朗日子流形，其极小 Maslov 指标为 $4$；
\end{enumerate}
% [源 B030]
则对任何系数域 $\Bbbk$，有结合代数的同构
\[
  HF^{\ast}(L,L;\Bbbk)\cong H^{\ast}(L;\Bbbk),
\]
\end{theorem*}

% [Tutor Note: B030 的结论是结合代数同构，而不只是向量空间同构：$HF^{\ast}(L,L;\Bbbk)$ 的 Floer 乘积对应 $H^{\ast}(L;\Bbbk)$ 的上同调环乘法（在 $HF^0$ 上即可核对）。]

% [Tutor Note: B027 原文 “Assuming on of the following condition is true” 中 “on” 应为 “one”（“condition” 亦应为复数）；两条假设分别针对“辛非球面且 $\pi_2(X,L)=0$”与“单调且极小 Maslov 指标 $4$”，彼此独立，不能合并成一条。]

% [源 B031]
特别地，当 $L$ 是拉格朗日环面且 $\Bbbk=\mathbb{R}$ 时，得到 $HF^{\ast}(L,L;\mathbb{R})\cong H^{\ast}(L;\mathbb{R})\cong\Lambda_{\mathbb{R}}\langle x_1,\dotsb,x_n\rangle$，因而在把基系数域换成 $\mathbb{C}$ 之后，它就接近于点对象。从几何上看，一种自然的扩域方式是给 $L$ 配上\emph{$U(1)$-局部系统}，等价地说，配上 $L$ 上的平凡复线束连同一个平坦 $U(1)$-联络。

% [Tutor Note: B031 的 “up to changing the base coefficient field to $\mathbb{C}$” 指基域扩张：$\Lambda_{\mathbb{R}}\langle x_1,\dotsb,x_n\rangle\otimes_{\mathbb{R}}\mathbb{C}\cong\Lambda_{\mathbb{C}}\langle x_1,\dotsb,x_n\rangle$；这与 B018 中 $k(P)$ 是 $k$ 的扩域的设定相呼应，也是把拉格朗日环面看成点对象的关键一步。]

% [源 B032]
这些讨论为 Fukaya 范畴中的点对象提供了一个自然的来源：彼此不相交的拉格朗日环面。然而，固定的辛流形中拉格朗日环面太多，并非它们都应当被算作“点”。为了把真正的“点对象”区分出来，Strominger--Yau--Zaslow [SYZ96] 建议我们应当关注\emph{特殊}的拉格朗日环面；在现代代数几何的语言中，这被期望等价于\emph{稳定}的概念 [TY02]。

% [Tutor Note: B032 首现 “special”（特殊的）与 “stable”（稳定的）两个说法，源文用 “expected to be an equivalent notion” 而非断言等价，译文写作“被期望等价于”以保留这一语气。]

% [源 B033]
我们知道确实可行的特殊情形之一是：$(X,\omega)$ 允许一个到某个光滑流形 $B$ 上的光滑纤维化，其纤维为拉格朗日环面 [Abo21]；这里假定 $\pi_2(B)=0$。一般地，我们只能期望 $X$ 的一个开稠密子集 $X^0$ 允许光滑的拉格朗日纤维化，而其补集由各种奇点构成。还有更糟糕的问题：对一般的辛流形 $(X,\omega)$，为使 Fukaya $A_{\infty}$-结构有定义，我们需要引入\emph{诺维科夫域}
\[
  \Lambda_{\Bbbk}:=\left\{\sum_{n=0}^{\infty}a_nT^{\alpha_n}\middle\vert a_n\in\Bbbk,\ \alpha_n\in\mathbb{R},\lim_{n\to\infty}\alpha_n=\infty\right\}
\]
来控制所计数的伪全纯圆盘的能量，以实现 Gromov 紧致性。因此，构造出的镜像空间自然地定义在系数域 $\Lambda_{\Bbbk}$ 上，而不是 $\Bbbk$ 上。Kontsevich 与 Soibelman [KS06] 建议，镜像应当是这一非阿基米德域 $\Lambda_{\Bbbk}$ 上的解析空间；把形式变量 $T$ 代入收敛的复数值，就归结为通常的复镜像空间。

% [Tutor Note: B033 的诺维科夫域要求 $\lim_{n\to\infty}\alpha_n=\infty$（即指数集有下界），于是对任何 $\alpha$ 只有有限多个 $n$ 满足 $\alpha_n\leq\alpha$，即能量有界的圆盘只涉及有限项，这正是 Gromov 紧致性所需；若把极限方向写反，级数将不构成域。$a_n\in\Bbbk$、$\alpha_n\in\mathbb{R}$ 中的 $n$ 只是可数枚举。]

% [源 B034]
拉格朗日环面纤维化的奇点可以分为若干不同的情形。注意，为使 $CF^{\ast}(L,L)$ 成为上链复形，需要证明 $(\mathfrak{m}^1)^2=0$；而在计算 $(\mathfrak{m}^1)^2$ 所用的伪全纯条带模空间的紧化中，会出现带单个输出的圆盘气泡。这个分层只在拉格朗日子流形 $L$ 围出 Maslov 指标 $\leq 2$ 的全纯圆盘时才会出现；因此，若这类圆盘全不出现，我们应当看到完美的镜像图景（如 [Abo21] 或 [PZ98] 中那样）。若存在 Maslov 指标为 $2$ 的全纯圆盘，则我们在 $CF^{\ast}(L,L)$ 的 $A_{\infty}$-结构中会得到一个额外的障碍项 $\mathfrak{m}^0$；但若假定 $L$ 是单调的，则 $\mathfrak{m}^0(L)=0$，从而 $HF^{\ast}(L,L)$ 仍可有定义。$\mathfrak{m}^0(L,\nabla)$ 在 $(L,\nabla)$ 的模空间上定义了一个光滑的全纯函数 [Aur07]，并充当镜像的\emph{Landau--Ginzburg 超势}。换言之，若假定 $(Y,J)$ 是镜像模空间、$W\colon Y\to\mathbb{C}$ 是对应的超势，则我们期望有范畴的等价
\[
  D^{\pi}\mathcal{Fuk}(X,\omega)\simeq HMF(Y,W),
\]
其中 $HMF$ 是 $(Y,W)$ 的\emph{矩阵分解}的同伦范畴，用来描述超势 $W$ 的奇点。

% [Tutor Note: B034 的 $\mathfrak{m}^0(L,\nabla)$ 依赖局部系统 $\nabla$：它计数 Maslov 指标 $2$ 的伪全纯圆盘，并按 $\nabla$ 沿边界的平行移动取迹；这解释了为何镜像一侧的模空间是 $(L,\nabla)$ 的对，而不是 $L$ 本身。]

% [源 B035]
$L$ 围出 Maslov 指标为 $0$ 的全纯圆盘的情形，在 [Aur07] 中也有初步研究。在良好的几何情形下，这类圆盘来自拉格朗日环面纤维化 $X\to\bar{B}$ 的奇点：其纤维是把某些子流形收缩成一点得到的奇异拉格朗日环面。此时我们看到：满足 $L$ 不围出任何 Maslov 指标 $0$ 圆盘的对 $(L,\nabla)$ 所构成的子空间是不连通的；而满足 $L$ 围出 Maslov 指标 $0$ 圆盘的对 $(L,\nabla)$ 所构成的空间形成实余维数为 $1$ 的子流形，称为该拉格朗日环面纤维化的\emph{墙}。在最一般的情形下，当奇异纤维是单收缩拉格朗日环面时，我们可以显式地理解单值化变换，并依据单值化信息把 $B$ 的不同连通分支粘起来。

% [Tutor Note: B035 的“墙”是实余维 $1$ 的子流形：$\dim_{\mathbb{R}}B=n$ 时它是 $(n-1)$ 维的；沿墙两侧 $(L,\nabla)$ 的模空间通过 Maslov 指标 $0$ 圆盘的贡献相互衔接，这正是后文“粘合模空间”的几何来源。]

% [源 B036]
然而，当我们希望保持某种对称性时，确实会出现非一般的情形，这要求我们对这种“模空间的粘合”有更好的理解。回忆对光滑的无障碍环面 $L$，$(L,\nabla)$ 的模空间局部地来自谱 $\operatorname{Spec}\operatorname{Sym}^{\ast}H^1(L,\nabla;\Lambda_{\mathbb{R}})$；而对奇异的 $L$，其拉格朗日 Floer 上链复形自然地有障碍，即 $\mathfrak{m}^0\neq 0$，从而 $(\mathfrak{m}^1)^2\neq 0$。此时我们可以引入闭链 $b\in CF^1(L,\nabla;\Lambda)$，使得 $\mathfrak{m}^1(b)=\mathfrak{m}^0$，以此作形变；这样的闭链称为\emph{约束上链}，形变后的 $A_{\infty}$-结构 $\mathfrak{m}^i_b$ 将满足所需的 $A_{\infty}$-关系，因此我们仍有良定义的 Floer 上同调。从几何上可以这样想象：把奇异的拉格朗日子流形 $L_0$ 形变为附近的光滑拉格朗日子流形 $L_{\eta}$，使其 Floer 上同调有良好定义。因此我们可以考虑所有这类形变构成的模空间，称为\emph{Maurer--Cartan 空间} $\mathcal{MC}(L_0)$，它应当在不同不连通的模空间之间起插值作用。

% [Tutor Note: B036 原文 “moduli spaces of paris” 中 “paris” 是 “pairs” 的笔误；又 $b\in CF^1(L,\nabla;\Lambda)$ 与 $\mathfrak{m}^1(b)=\mathfrak{m}^0$ 自洽：在 $|\mathfrak{m}^k|=1-k$ 的约定下 $\mathfrak{m}^0$ 的次数为 $1$，故约束上链须取在次数 $1$ 处。]

% [源 B037]
至此，关于 SYZ 镜像对称构造如何运作并给出我们想要的镜像空间，粗略的理论框架就讲完了。然而，正如听众将在本次报告中看到的，仍有大量奥秘有待理解。在随后的报告中，我们将更详细地看到一些计算，以及若干关于 SYZ 与同调镜像对称的严格陈述。
